\documentclass[12pt]{article}

\usepackage[T2A]{fontenc}
\usepackage[utf8]{inputenc}
\usepackage[russian,english]{babel}
\usepackage{array}
\usepackage{enumitem}
\usepackage{geometry}
\usepackage{graphicx}
\usepackage{hyperref}
\usepackage{makecell}
\usepackage{tabularx}
\usepackage{tikz}
\usepackage{xltabular}
\usepackage[table]{xcolor}

\geometry{
    a4paper,
    left=1.5cm,
    right=1.5cm,
    top=1.5cm,
    bottom=2.5cm
}

\hypersetup{hidelinks}

\setlength{\parindent}{0cm}
\setlength{\parskip}{0.35\baselineskip}
\setlength{\emergencystretch}{1.5em}
\setlength{\tabcolsep}{7pt}
\renewcommand{\arraystretch}{1.12}
\setcellgapes{2pt}
\makegapedcells
\keepXColumns

\definecolor{commitred}{RGB}{236, 48, 48}
\definecolor{commitblue}{RGB}{52, 83, 255}
\newcolumntype{L}[1]{>{\raggedright\arraybackslash}p{#1}}
\newcolumntype{Y}{>{\raggedright\arraybackslash}X}
\newcommand{\cmdcell}[1]{\parbox[t]{\linewidth}{\ttfamily\small\raggedright\strut #1\strut}}
\newcommand{\commitnode}[4]{%
    \fill[#3] (#1,#2) circle (4.3pt);
    \node[font=\bfseries\scriptsize] at (#1,#2 + 0.42) {#4};
}

\begin{document}

\pagenumbering{gobble}
\begin{titlepage}
    \centering
    Федеральное государственное автономное \\
    образовательное учреждение высшего образования \\
    \textbf{<<Национальный исследовательский университет ИТМО>>} \\

    \vfill
    Основы программной инженерии \\
    \vspace{0.5\baselineskip}
    \textbf{Лабораторная работа №2} \\
    \vspace{0.5\baselineskip}
    Вариант 1400 \\

    \vfill

    \flushright
    Выполнили:\\
    Чупров Иван Андреевич P3214 ОПИ 2.2\\
    Лесников Владимир Алексеевич P3214 ОПИ 2.2\\

    \vspace{4\baselineskip}

    \centering
    г. Санкт-Петербург, 2025

\end{titlepage}

\clearpage
\pagenumbering{arabic}
\tableofcontents
\newpage

\section{Задание}

\subsection{Цель работы}

Сконфигурировать в домашнем каталоге пользователя репозитории
\texttt{svn} и \texttt{git}, загрузить в них начальную ревизию файлов
с исходными кодами в соответствии с вариантом и воспроизвести
последовательность операций над исходным кодом по выданной блок-схеме.

При составлении последовательности команд учитываются следующие условия:

\begin{itemize}[leftmargin=1.5em]
    \item цвет элементов схемы указывает на пользователя,
    совершившего действие: красный --- первый, синий --- второй;
    \item цифры над узлами обозначают номера ревизий;
    \item при слияниях необходимо разрешать возникающие конфликты.
\end{itemize}

\subsection{Блок-схема варианта}

\begin{figure}[h]
    \centering
    \begin{tikzpicture}[x=0.62cm, y=0.78cm, line width=0.9pt]
        \draw[commitred] (0,2) -- (2,2) -- (7,2) -- (12.5,2) -- (13.5,2);
        \draw[commitred] (0,2) -- (0,0.55) -- (1,0.55) -- (8,0.55) -- (9,0.55) -- (11.55,0.55) -- (12.5,0.55) -- (12.5,2);
        \draw[commitblue] (2,2) -- (2,1.15) -- (3,1.15) -- (4,1.15) -- (5,1.15) -- (6,1.15) -- (10,1.15) -- (11,1.15) -- (11.55,1.15) -- (11.55,0.55);

        \commitnode{0}{2}{commitred}{r0}
        \commitnode{1}{0.55}{commitred}{r1}
        \commitnode{2}{2}{commitred}{r2}
        \commitnode{3}{1.15}{commitblue}{r3}
        \commitnode{4}{1.15}{commitblue}{r4}
        \commitnode{5}{1.15}{commitblue}{r5}
        \commitnode{6}{1.15}{commitblue}{r6}
        \commitnode{7}{2}{commitred}{r7}
        \commitnode{8}{0.55}{commitred}{r8}
        \commitnode{9}{0.55}{commitred}{r9}
        \commitnode{10}{1.15}{commitblue}{r10}
        \commitnode{11}{1.15}{commitblue}{r11}
        \commitnode{11.55}{0.55}{commitred}{r12}
        \commitnode{12.5}{2}{commitred}{r13}
        \commitnode{13.5}{2}{commitred}{r14}
    \end{tikzpicture}
    \caption{Блок-схема варианта 1400}
\end{figure}

\section{Ход работы}

Для краткости в таблицах ниже используется обозначение
\texttt{apply\_snapshot N}. Под ним понимается замена содержимого рабочего
каталога состоянием из архива \texttt{commitN.zip}.
Для \texttt{git} после этого выполняется \texttt{git add -A}, для
\texttt{svn} --- \texttt{svn add --force .} после удаления старого состояния.

\subsection{Создание и конфигурация репозитория Git}

\noindent
\begin{tabularx}{\textwidth}{|L{0.18\textwidth}|Y|}
\hline
\rowcolor{white}
\textbf{Этап} & \textbf{Команды} \\
\hline
Подготовка &
\cmdcell{cd ~\\
mkdir git-usage\\
cd git-usage\\
git init\\
git config user.name "red"\\
git config user.email "red@example.com"\\
git config merge.conflictstyle diff3} \\
\hline
\end{tabularx}

\subsection{Последовательность ревизий Git}

\begin{xltabular}{\textwidth}{|L{0.10\textwidth}|>{\hsize=1.45\hsize\raggedright\arraybackslash}X|>{\hsize=0.75\hsize\raggedright\arraybackslash}X|}
\hline
\rowcolor{white}
\textbf{Коммит} & \textbf{Команды} & \textbf{Комментарий} \\
\hline
\endfirsthead
\hline
\rowcolor{white}
\textbf{Коммит} & \textbf{Команды} & \textbf{Комментарий} \\
\hline
\endhead
r0 &
\cmdcell{apply\_snapshot 0\\
git commit -m "r0"} &
Начальная ревизия. \\
\hline
r1 &
\cmdcell{git checkout -b red-bottom\\
apply\_snapshot 1\\
git commit -m "r1"} &
Создание нижней красной ветки. \\
\hline
r2 &
\cmdcell{git checkout main\\
apply\_snapshot 2\\
git commit -m "r2"} &
Продолжение верхней красной ветки. \\
\hline
r3 &
\cmdcell{git config user.name "blue"\\
git config user.email "blue@example.com"\\
git checkout -b blue\\
apply\_snapshot 3\\
git commit -m "r3"} &
Начало синей ветки от \texttt{r2}. \\
\hline
r4 &
\cmdcell{apply\_snapshot 4\\
git commit -m "r4"} &
Следующая ревизия синей ветки. \\
\hline
r5 &
\cmdcell{apply\_snapshot 5\\
git commit -m "r5"} &
Продолжение синей ветки. \\
\hline
r6 &
\cmdcell{apply\_snapshot 6\\
git commit -m "r6"} &
Продолжение синей ветки. \\
\hline
r7 &
\cmdcell{git config user.name "red"\\
git config user.email "red@example.com"\\
git checkout main\\
apply\_snapshot 7\\
git commit -m "r7"} &
Возврат в верхнюю красную ветку. \\
\hline
r8 &
\cmdcell{git checkout red-bottom\\
apply\_snapshot 8\\
git commit -m "r8"} &
Продолжение нижней красной ветки. \\
\hline
r9 &
\cmdcell{apply\_snapshot 9\\
git commit -m "r9"} &
Продолжение нижней красной ветки. \\
\hline
r10 &
\cmdcell{git config user.name "blue"\\
git config user.email "blue@example.com"\\
git checkout blue\\
apply\_snapshot 10\\
git commit -m "r10"} &
Возврат в синюю ветку. \\
\hline
r11 &
\cmdcell{apply\_snapshot 11\\
git commit -m "r11"} &
Подготовка к первому слиянию. \\
\hline
r12 &
\cmdcell{git config user.name "red"\\
git config user.email "red@example.com"\\
git checkout red-bottom\\
git merge --no-ff --no-commit blue\\
apply\_snapshot 12\\
git commit -m "r12"} &
Слияние \texttt{blue -> red-bottom}, конфликтующие файлы приводятся к состоянию \texttt{commit12.zip}. \\
\hline
r13 &
\cmdcell{git checkout main\\
git merge --no-ff --no-commit red-bottom\\
apply\_snapshot 13\\
git commit -m "r13"} &
Слияние \texttt{red-bottom -> main}, конфликтующие файлы приводятся к состоянию \texttt{commit13.zip}. \\
\hline
r14 &
\cmdcell{apply\_snapshot 14\\
git commit -m "r14"} &
Финальная ревизия. \\
\hline
\end{xltabular}

\subsection{Создание и конфигурация репозитория SVN}

\noindent
\begin{tabularx}{\textwidth}{|L{0.18\textwidth}|Y|}
\hline
\rowcolor{white}
\textbf{Этап} & \textbf{Команды} \\
\hline
Подготовка &
\cmdcell{cd ~\\
mkdir svn-usage\\
cd svn-usage\\
svnadmin create repo\\
REPO\_URL="file://\$HOME/svn-usage/repo"\\
svn mkdir "\$REPO\_URL/trunk" "\$REPO\_URL/branches"\\
-m "Init repository layout"\\
svn checkout "\$REPO\_URL/trunk" workDir\\
cd workDir} \\
\hline
\end{tabularx}

\subsection{Последовательность ревизий SVN}

Команды \texttt{svn mkdir} и \texttt{svn copy} создают служебные ревизии
Subversion, поэтому их внутренние номера не совпадают с номерами узлов
на блок-схеме. В таблице приведены действия, соответствующие именно
ревизиям варианта.

\begin{xltabular}{\textwidth}{|L{0.10\textwidth}|>{\hsize=1.45\hsize\raggedright\arraybackslash}X|>{\hsize=0.75\hsize\raggedright\arraybackslash}X|}
\hline
\rowcolor{white}
\textbf{Коммит} & \textbf{Команды} & \textbf{Комментарий} \\
\hline
\endfirsthead
\hline
\rowcolor{white}
\textbf{Коммит} & \textbf{Команды} & \textbf{Комментарий} \\
\hline
\endhead
r0 &
\cmdcell{apply\_snapshot 0\\
svn commit -m "r0" --username red} &
Начальная ревизия в \texttt{trunk}. \\
\hline
r1 &
\cmdcell{svn copy "\$REPO\_URL/trunk"\\
"\$REPO\_URL/branches/red-bottom"\\
-m "Create red-bottom branch" --username red\\
svn switch "\$REPO\_URL/branches/red-bottom"\\
apply\_snapshot 1\\
svn commit -m "r1" --username red} &
Создание нижней красной ветки. \\
\hline
r2 &
\cmdcell{svn switch "\$REPO\_URL/trunk"\\
apply\_snapshot 2\\
svn commit -m "r2" --username red} &
Продолжение верхней красной ветки. \\
\hline
r3 &
\cmdcell{svn copy "\$REPO\_URL/trunk"\\
"\$REPO\_URL/branches/blue"\\
-m "Create blue branch" --username red\\
svn switch "\$REPO\_URL/branches/blue"\\
apply\_snapshot 3\\
svn commit -m "r3" --username blue} &
Начало синей ветки от \texttt{r2}. \\
\hline
r4 &
\cmdcell{apply\_snapshot 4\\
svn commit -m "r4" --username blue} &
Следующая ревизия синей ветки. \\
\hline
r5 &
\cmdcell{apply\_snapshot 5\\
svn commit -m "r5" --username blue} &
Продолжение синей ветки. \\
\hline
r6 &
\cmdcell{apply\_snapshot 6\\
svn commit -m "r6" --username blue} &
Продолжение синей ветки. \\
\hline
r7 &
\cmdcell{svn switch "\$REPO\_URL/trunk"\\
apply\_snapshot 7\\
svn commit -m "r7" --username red} &
Возврат в верхнюю красную ветку. \\
\hline
r8 &
\cmdcell{svn switch "\$REPO\_URL/branches/red-bottom"\\
apply\_snapshot 8\\
svn commit -m "r8" --username red} &
Продолжение нижней красной ветки. \\
\hline
r9 &
\cmdcell{apply\_snapshot 9\\
svn commit -m "r9" --username red} &
Продолжение нижней красной ветки. \\
\hline
r10 &
\cmdcell{svn switch "\$REPO\_URL/branches/blue"\\
apply\_snapshot 10\\
svn commit -m "r10" --username blue} &
Возврат в синюю ветку. \\
\hline
r11 &
\cmdcell{apply\_snapshot 11\\
svn commit -m "r11" --username blue} &
Подготовка к первому слиянию. \\
\hline
r12 &
\cmdcell{svn switch "\$REPO\_URL/branches/red-bottom"\\
svn merge "\$REPO\_URL/branches/blue" --accept postpone\\
svn resolve --accept working -R .\\
apply\_snapshot 12\\
svn resolve --accept working -R .\\
svn commit -m "r12" --username red} &
Слияние \texttt{blue -> red-bottom}, конфликтующие файлы приводятся к состоянию \texttt{commit12.zip}. \\
\hline
r13 &
\cmdcell{svn switch "\$REPO\_URL/trunk"\\
svn merge "\$REPO\_URL/branches/red-bottom" --accept postpone\\
svn resolve --accept working -R .\\
apply\_snapshot 13\\
svn resolve --accept working -R .\\
svn commit -m "r13" --username red} &
Слияние \texttt{red-bottom -> trunk}, конфликтующие файлы приводятся к состоянию \texttt{commit13.zip}. \\
\hline
r14 &
\cmdcell{apply\_snapshot 14\\
svn commit -m "r14" --username red} &
Финальная ревизия. \\
\hline
\end{xltabular}

\subsection{Разрешение конфликтов}

Конфликты возникали в точках слияния \texttt{r12} и \texttt{r13}, когда в
объединяемых ветках были изменены одни и те же файлы. В обоих случаях
разрешение выполнялось одинаково:

\begin{itemize}[leftmargin=1.5em]
    \item запускалось слияние нужной ветки;
    \item рабочее дерево приводилось к целевому состоянию из архивов
    \texttt{commit12.zip} и \texttt{commit13.zip};
    \item результат помечался как разрешённый и фиксировался коммитом:
    \texttt{git add -A} для Git и
    \texttt{svn resolve --accept working -R .} для SVN.
\end{itemize}

Такой подход гарантирует, что после merge содержимое файлов точно совпадает
с состоянием, заданным вариантом.

\section{Выводы}

В ходе работы были воспроизведены одинаковые состояния проекта в
\texttt{git} и \texttt{svn} в соответствии с вариантом 1400. Были
отработаны создание и конфигурация репозиториев, работа с ветками,
переключение пользователей и разрешение конфликтов при слиянии.
Практика показала различия между распределённой и централизованной
моделями контроля версий и закрепила базовые команды обеих систем.

\end{document}
